% =====================================================================
% Appendix H --- The model: kernel, labour closures, financing closures
% ---------------------------------------------------------------------
% Job: the formal detail that Section 4 deliberately leaves out. Section 4
%      keeps only the input--output table and the two optimization problems;
%      the functional forms of their solutions, the five labour closures and
%      the three financing closures live here.
% Moved from Section 4 (2026-10-12); equation labels kept unchanged so the
% \eqref's in Sections 3--6 and the tables still resolve.
% Formulations: registry/closures.toml; ADR-0002, ADR-0019, ADR-0020,
% ADR-0022. The full stacked system and the proofs are in Appendix C
% (app:eqkernel); the A-bill accounting is in Appendix D (app:kernel).
% =====================================================================
\section{The model: kernel and closure formulations}\label{app:model}

\noindent This appendix states the formal model behind
Section~\ref{sec:kernel}: the functional forms that solve the producer and
household problems of Section~\ref{subsec:optimization}, the five labour
closures and the three financing closures. The stacked equation system used
in the proofs is in Appendix~\ref{app:eqkernel}; the intermediate-bill and
external-account accounting is in Appendix~\ref{app:kernel}.

% Appendix equations share the A-prefixed counter with the other appendices; keep the
% PDF anchors distinct.
\begingroup
\renewcommand*{\theHequation}{model.\arabic{equation}}

\subsection{The kernel}\label{app:model-kernel}

Subscript $0$ denotes the baseline. Each sector's technology is a nested CES
with constant returns: labour and an intermediate composite substitute with
elasticity $\epsilon$, and intermediate inputs substitute with elasticity
$\theta$ inside the composite. Cost minimization gives the intermediate
price index and the unit cost
\begin{align}
P_i(p) &= \Big(\textstyle\sum_{u=1}^N \Omega_{iu}\,p_u^{1-\theta}\Big)^{1/(1-\theta)},
\label{eq:model-ip}\\
c_i(p,w_i;A_i) &= \Big(A_i^{\epsilon-1}\big(\alpha_i\,w_i^{1-\epsilon}
  + (1-\alpha_i)\,P_i^{1-\epsilon}\big)\Big)^{1/(1-\epsilon)},
\label{eq:model-cost}
\end{align}
with the Cobb--Douglas limits at $\theta=1$ and $\epsilon=1$. Because
$\sum_u\Omega_{iu}=1$, the unit cost equals one at the calibration point
$(p,w,A)=(\one,\one,\one)$. Zero profit requires $p_i = c_i(p,w_i;A_i)$.
Shephard's lemma, evaluated at zero profit, gives the cost-minimizing
labour demand
\begin{equation}
L_i^{\mathrm{cm}}(p,y_i,w_i)
 = \alpha_i A_i^{\epsilon-1} y_i
   \left(\frac{p_i}{w_i}\right)^{\epsilon},
\label{eq:closure-labour-demand}
\end{equation}
where $y_i$ is gross output, $A_i$ productivity and $\alpha_i$ the calibrated
labour share. Intermediate demand follows in the same way from the
composite~\eqref{eq:model-ip} and is charged with the domestic bill
(Appendix~\ref{app:kernel}; explicit form in~\eqref{eq:int}).

The household has CES preferences with elasticity $\sigma$ and calibrated
weights $\omega_i$. The consumption price index, which is also the
numeraire, is
\begin{equation}
  \Pi(p) = \left(\sum_{i=1}^N \omega_i p_i^{1-\sigma}\right)^{1/(1-\sigma)},
  \qquad \sum_{i=1}^N \omega_i = 1,
  \label{eq:labour-cpi}
\end{equation}
with the Cobb--Douglas limit at $\sigma = 1$. Given household consumption
expenditure $E^h$, utility maximization gives
\begin{equation}
c_i^h = E^h\frac{\omega_i p_i^{-\sigma}}{\sum_j\omega_j p_j^{1-\sigma}}.
\label{eq:model-household}
\end{equation}
Goods-market clearing for each of the $N$ sectors reads
\begin{equation}
y_i = \mathit{int}_i + (1-m_i)\big(c_i^h + g_i^G + \mathit{inv}_i + g_i\big) + X_i,
\label{eq:model-clearing}
\end{equation}
where $m_i$ is the sectoral import margin, $g^G$ baseline government
purchases, $\mathit{inv}$ non-programme investment, $X$ exports and $g$ the
programme bundle (zero under F1). Every closure below enforces all $N$
clearings~\eqref{eq:model-clearing}.

\subsection{The five labour closures}\label{app:labour-closures}

Let $L_i$ denote employment in sector $i$, $\bar{L}_i$ its baseline value,
and $\Lbar = \sum_{i=1}^N \bar{L}_i$ total baseline employment. Sectoral
wages are $w_i$; with mobile labour they equal one common wage $w$. All
closures use the same labour demand~\eqref{eq:closure-labour-demand}; they
differ in the wage and labour-supply restrictions imposed on it.

\paragraph{The no-reallocation closure: immobile labour.}
The no-reallocation row fixes employment in every sector and solves a
sector-specific wage to clear each labour market:
\begin{equation}
L_i = L_i^{\mathrm{cm}}(p,y_i,w_i) = \bar{L}_i,
\qquad i=1,\ldots,N, \qquad \eta = 0.
\label{eq:closure-bf}
\end{equation}
Thus $\sum_i L_i = \Lbar$, but wages may differ across sectors.
The allocation parameter $\eta$ selects only the endpoints $\{0,1\}$;
its mobile endpoint $\eta = 1$ is the standard CGE closure below, not an
additional closure.

\paragraph{The standard CGE closure: mobile labour at full employment.}
Labour reallocates across sectors at one flexible wage, while total
employment remains fixed:
\begin{equation}
L_i = L_i^{\mathrm{cm}}(p,y_i,w),
\qquad \sum_{i=1}^N L_i = \Lbar,
\qquad \eta = 1.
\label{eq:closure-alpha}
\end{equation}

\paragraph{The elastic labour-supply closure: sectoral supply curves.}
We use the uniform-elasticity sectoral form: each sector has its own wage
and supply curve, with one common labour-supply elasticity $\eta_s$:
\begin{equation}
L_i = L_i^{\mathrm{cm}}(p,y_i,w_i)
 = \bar{L}_i \left[\frac{w_i/\Pi(p)}{w_{i0}/\Pi_0}\right]^{\eta_s},
\qquad \eta_s \geq 0, \qquad i=1,\ldots,N.
\label{eq:closure-beta}
\end{equation}
Equal supply elasticities do not impose equal wages: $N$ separate labour
markets determine $w_i$. Here $\eta_s$ is a labour-supply elasticity, not
the allocation parameter $\eta$; setting $\eta_s = 0$ gives exactly the
no-reallocation row. The calibrated real-wage anchors are $w_{i0}/\Pi_0 = 1$.
The deflator remains the common consumption price index, not the sector's
own goods price $p_i$. The supply curves are reduced forms without a
labour--leisure income effect. The headline panels use $\eta_s = 0.5$;
the elasticity ladder and heterogeneous sectoral variants are reported in
Appendix~\ref{app:variants}. The single-wage form of this closure is
retained there only as a distinct benchmark; its zero-elasticity limit is
the standard CGE closure, not the no-reallocation one.

\paragraph{The fixed real wage closure.}
The real wage is pegged at its baseline value, while employment is
endogenous and uncapped:
\begin{equation}
\frac{w}{\Pi(p)} = \bar{w} = 1,
\qquad L_i = L_i^{\mathrm{cm}}(p,y_i,w),
\qquad \eta = 1.
\label{eq:closure-gamma}
\end{equation}
Unlike the standard CGE closure, this closure does not impose
$\sum_i L_i = \Lbar$.
It is a two-sided wage peg, not a one-sided wage floor with unemployment
rationing. On the demand-only baseline branch, $\Pi(p)=1$ and the
implemented wage pin is $w=1$.

\paragraph{The Leontief closure: the fixed-coefficient corner.}
The Leontief closure keeps the wage and employment rule of the fixed real
wage closure while taking the
substitution elasticities to their Leontief limits:
\begin{equation}
\frac{w}{\Pi(p)} = \bar{w} = 1,
\qquad \eta = 1,
\qquad (\theta,\epsilon,\sigma) \longrightarrow (0^+,0^+,0^+).
\label{eq:closure-delta}
\end{equation}
The Leontief closure is therefore a corner of the fixed real wage closure
with fixed-coefficient technology, rather than an independent
labour-market mechanism.

\subsection{The three financing closures}\label{app:financing-closures}
% Programme-scale F1 tilt: experiments/run.jl, f1_tilt_weights.

Each labour closure is crossed with three ways of accommodating the same
programme scale and sectoral incidence. Let $g_i = kG_0\psi_i$ denote the
reference programme bundle, where $G_0$ is its baseline-price value,
$k \geq 0$ scales the programme, and $\psi_i \geq 0$ with
$\sum_i\psi_i = 1$ describes its sectoral incidence. F2 and F3 add this
bundle to final demand; F1 instead uses its scale and incidence to redirect
household spending.

Baseline government purchases $g_i^G$ remain tax-financed in all three
cases. Define the wage bill and their current-price tax cost as
\begin{equation}
W = \sum_{i=1}^N w_i L_i,
\qquad T_G(p) = \sum_{i=1}^N p_i g_i^G.
\label{eq:financing-baseline-budget}
\end{equation}
With a common wage, $W = w\sum_i L_i$. Let $E$ denote after-tax household
resources, including the net external transfer $R$ from the rest of the
world,\footnote{$R$ is the variable called $F$ in the \texttt{OpenCGE}
kernel and in the run manifests. The labels F1--F3 name financing closures,
not a variable.} and let $s$ be the
calibrated saving rate. Household consumption expenditure is
$E^h = (1-s)E$. Gross household purchases $c_i^h$ exhaust this budget:
\begin{equation}
\sum_{i=1}^N p_i c_i^h = E^h,
\qquad c_i^{\mathrm{dom}} = (1-m_i)c_i^h.
\label{eq:financing-household-budget}
\end{equation}
Only the domestic content enters domestic goods-market
clearing~\eqref{eq:model-clearing}; likewise, the programme contributes
$(1-m_i)g_i$ under F2 and F3.

\paragraph{F1: budget-neutral preference reallocation.}
F1 changes the household's CES weights without adding a public-demand
bundle. For positive preference shifters $d_i$, the weights and gross
household demand become
\begin{align}
\widetilde{\omega}_i
 &= \frac{\omega_i d_i}{\sum_j\omega_j d_j},
\qquad \sum_i\widetilde{\omega}_i = 1,
\label{eq:closure-f1-weights}\\
c_i^h
 &= E^h\frac{\widetilde{\omega}_i p_i^{-\sigma}}
 {\sum_j\widetilde{\omega}_j p_j^{1-\sigma}},
\qquad E = W-T_G(p)+R.
\label{eq:closure-f1}
\end{align}
The programme-scale tilt is $d_i = 1+kG_0\psi_i^+/c_{i0}^h$ in sectors
with positive baseline household purchases, and $d_i=1$ otherwise;
$\psi^+$ is the programme incidence restricted to those sectors and
renormalized to sum to one. The CES normalizer preserves
\eqref{eq:financing-household-budget} at every price vector. Thus increased
spending shares in programme sectors are paid for by reduced shares
elsewhere, conditional on the household budget; equilibrium income need
not remain unchanged. F1 is a demand-composition experiment, not an
autonomous investment injection or multiplier.

\paragraph{F2: tax-financed public investment.}
F2 adds the real programme bundle $g$ and levies an additional household
tax equal to its value at equilibrium prices:
\begin{equation}
T(p) = \sum_{i=1}^N p_i g_i,
\qquad E = W-T_G(p)-T(p)+R,
\qquad B_{\mathrm{gov}}=0.
\label{eq:closure-f2}
\end{equation}
Total domestic tax revenue $T_G(p)+T(p)$ therefore balances the purchases
$g^G+g$, priced at current rather than baseline prices. The tax is
implemented as the endogenous wage-income rate
$\tau=[T_G(p)+T(p)]/W$; with one representative household its burden is
effectively lump-sum. The programme has an explicit domestic budget
counterparty and reduces household resources directly by $T(p)$.

\paragraph{F3: externally financed public investment.}
F3 adds the same real bundle but imposes no additional domestic programme
tax. Its current-price cost is booked as external programme financing:
\begin{equation}
B_{\mathrm{gov}} = \sum_{i=1}^N p_i g_i,
\qquad E = W-T_G(p)+R.
\label{eq:closure-f3}
\end{equation}
The booking $B_{\mathrm{gov}}$ is distinct from the net external transfer
$R$ entering household resources after tax. The latter is solved
endogenously in the no-reallocation, the standard CGE and both
elastic-supply forms, but is fixed at zero in the two fixed-wage closures,
the fixed real wage closure and its Leontief corner. With no unfinanced
demand injection, all closures clear
all $N$ goods markets and satisfy the external-account identity
\begin{equation}
S + T_{\mathrm{int}} + M - (I+X) = R+B_{\mathrm{gov}},
\qquad S=sE,
\label{eq:financing-external-account}
\end{equation}
where $I$ is non-programme investment, $X$ exports, $M$ total imports
(including intermediate and programme imports), and $T_{\mathrm{int}}$
the separately booked product-tax leak on intermediate use; all are
current-price values. Here $B_{\mathrm{gov}}=0$ under F1 and F2, so the net
external position is $R+B_{\mathrm{gov}}$, not the programme cost alone.
F3 is an accounting open-economy closure, not a monetary mechanism: the
static model cannot distinguish external borrowing from a transfer and
does not model interest payments, repayment or exchange-rate adjustment.
\endgroup
